%% ============================================================
%%  Chapter 4 — Authentication and Role-Based Access Control
%%  04_authentication.tex
%% ============================================================
\chapter{Authentication and Role-Based Access Control}
\label{chap:auth}

\chapterintro{%
  This chapter provides a complete technical specification of the
  authentication and authorisation subsystem. It covers the Laravel
  Passport OAuth2 implementation in depth---including the token
  issuance pipeline, the middleware chain, and the revocation
  strategy---followed by a precise definition of the three-tier role
  hierarchy, a comprehensive mapping of every API endpoint to its
  required role, and a formal analysis of the mechanisms by which the
  system prevents both horizontal and vertical privilege escalation.
}

%% ─────────────────────────────────────────────────────────────────────────
\section{Authentication Architecture Overview}
\label{sec:auth-overview}
%% ─────────────────────────────────────────────────────────────────────────

The Clinic Management API uses a \textbf{stateless, token-based
authentication model} built on the OAuth2 specification, implemented
through the Laravel Passport package. Every request that accesses a
protected resource must present a cryptographically signed Bearer token
in the HTTP \texttt{Authorization} header. The server neither sets
cookies nor maintains any session state between requests; the token
itself is the sole proof of identity.

The choice of OAuth2 Personal Access Tokens over a simpler shared-secret
or JWT-only scheme is deliberate. Passport's token infrastructure stores
a record of every issued token in the \texttt{oauth\_access\_tokens}
table, which enables \textbf{server-side revocation}: when a user logs
out, the token is invalidated at the database level. Any subsequent
request bearing that token will be rejected with a \stunauth{} response,
regardless of whether the token's cryptographic signature is still valid.
This property is essential in a clinical environment where a lost or
stolen device must be immediately locked out.

\begin{figure}[h!]
  \centering
  \begin{tcolorbox}[
    enhanced, arc=4pt, boxrule=1pt,
    colback=LightGray, colframe=RuleGray,
    left=12pt, right=12pt, top=10pt, bottom=10pt,
    width=\linewidth,
  ]
  \centering
  \textbf{Client Request}
  \quad\textcolor{AccentTeal}{$\longrightarrow$}\quad
  \textcolor{PrimaryBlue}{\textbf{auth:api Middleware}}
  \quad\textcolor{AccentTeal}{$\longrightarrow$}\quad
  \textcolor{PrimaryBlue}{\textbf{RoleMiddleware}}
  \quad\textcolor{AccentTeal}{$\longrightarrow$}\quad
  \textcolor{PrimaryBlue}{\textbf{ClinicScopeMiddleware}}
  \quad\textcolor{AccentTeal}{$\longrightarrow$}\quad
  \textcolor{MethodPost}{\textbf{Controller}}

  \vspace{6pt}
  \small
  Each middleware either advances the request or terminates it with an
  appropriate HTTP error response.
  \end{tcolorbox}
  \caption{The three-stage middleware pipeline every protected request
    traverses before reaching a controller.}
  \label{fig:middleware-pipeline}
\end{figure}

%% ─────────────────────────────────────────────────────────────────────────
\section{Laravel Passport: Installation and Configuration}
\label{sec:passport-setup}
%% ─────────────────────────────────────────────────────────────────────────

\subsection{Package Installation and Key Generation}

Laravel Passport is installed as a Composer dependency. After
installation, the \texttt{passport:install} Artisan command generates
the RSA key pair that Passport uses to sign and verify access tokens.
The private key signs outgoing tokens; the public key verifies incoming
ones. These keys must be present in the deployment environment and must
never be committed to version control.

\begin{lstlisting}[style=bash, caption={Installing Passport and generating
  encryption keys.}]
composer require laravel/passport

php artisan passport:install
# Generates storage/oauth-private.key and storage/oauth-public.key
# Runs the oauth_* migrations (clients, tokens, auth codes)

php artisan migrate
\end{lstlisting}

\subsection{Model Integration}

The \phpclass{User} Eloquent model uses the
\texttt{Laravel\textbackslash{}Passport\textbackslash{}HasApiTokens}
trait, which adds the \texttt{createToken()} method and the
\texttt{tokens} relationship used throughout the authentication service.

\begin{lstlisting}[style=php, caption={\texttt{app/Models/User.php} ---
  Passport trait integration and role enum casting.}]
<?php

namespace App\Models;

use App\Enums\UserRole;
use Illuminate\Database\Eloquent\Concerns\HasUuids;
use Illuminate\Foundation\Auth\User as Authenticatable;
use Illuminate\Database\Eloquent\Factories\HasFactory;
use Illuminate\Database\Eloquent\Relations\HasMany;
use Laravel\Passport\HasApiTokens;

class User extends Authenticatable
{
    use HasApiTokens, HasFactory, HasUuids;

    protected $fillable = [
        'name',
        'email',
        'password',
        'role',
        'clinic_id',
        'is_active',
    ];

    protected $hidden = [
        'password',
        'remember_token',
    ];

    protected function casts(): array
    {
        return [
            'role'      => UserRole::class,
            'is_active' => 'boolean',
            'password'  => 'hashed',
        ];
    }

    public function clinic(): BelongsTo
    {
        return $this->belongsTo(Clinic::class);
    }

    public function appointments(): HasMany
    {
        return $this->hasMany(Appointment::class, 'doctor_id');
    }

    public function schedules(): HasMany
    {
        return $this->hasMany(DoctorSchedule::class, 'doctor_id');
    }
}
\end{lstlisting}

\subsection{Passport Registration in the Application Bootstrap}

In Laravel 12, Passport is registered in
\texttt{bootstrap/app.php} rather than a dedicated service provider.
The \texttt{withRouting()} call integrates Passport's own OAuth2
endpoint group, and \texttt{withMiddleware()} registers the
\texttt{auth:api} alias.

\begin{lstlisting}[style=php, caption={\texttt{bootstrap/app.php} ---
  Passport registration and API guard configuration.}]
<?php

use Illuminate\Foundation\Application;
use Illuminate\Foundation\Configuration\Exceptions;
use Illuminate\Foundation\Configuration\Middleware;
use Laravel\Passport\Http\Middleware\CheckForAnyScope;
use App\Http\Middleware\RoleMiddleware;
use App\Http\Middleware\ClinicScopeMiddleware;

return Application::configure(basePath: dirname(__DIR__))
    ->withRouting(
        api: __DIR__.'/../routes/api.php',
        apiPrefix: 'api',
        commands: __DIR__.'/../routes/console.php',
        health: '/up',
    )
    ->withMiddleware(function (Middleware $middleware) {
        $middleware->alias([
            'role'         => RoleMiddleware::class,
            'clinic.scope' => ClinicScopeMiddleware::class,
        ]);
    })
    ->withExceptions(function (Exceptions $exceptions) {
        // Exception handler bindings — see Chapter 6.
    })->create();
\end{lstlisting}

\subsection{API Guard Configuration}

Laravel's \texttt{config/auth.php} designates Passport as the driver
for the \texttt{api} guard. The \texttt{provider} points to the
\texttt{users} provider, which resolves to the \phpclass{User} model.

\begin{lstlisting}[style=php, caption={\texttt{config/auth.php} ---
  API guard using the Passport driver.}]
'guards' => [
    'web' => [
        'driver'   => 'session',
        'provider' => 'users',
    ],

    'api' => [
        'driver'   => 'passport',
        'provider' => 'users',
    ],
],

'providers' => [
    'users' => [
        'driver' => 'eloquent',
        'model'  => App\Models\User::class,
    ],
],
\end{lstlisting}

%% ─────────────────────────────────────────────────────────────────────────
\section{The Authentication Service and Token Lifecycle}
\label{sec:auth-service}
%% ─────────────────────────────────────────────────────────────────────────

\subsection{The UserRole Enum}

Before examining the authentication service, it is necessary to
understand the \phpclass{UserRole} enum that is central to all
role-based decisions throughout the application.

\begin{lstlisting}[style=php, caption={\texttt{app/Enums/UserRole.php} ---
  PHP 8.1 backed string enum for the three system roles.}]
<?php

namespace App\Enums;

enum UserRole: string
{
    case SUPER_ADMIN = 'super_admin';
    case DOCTOR      = 'doctor';
    case ASSISTANT   = 'assistant';

    public function label(): string
    {
        return match($this) {
            UserRole::SUPER_ADMIN => 'Super Administrator',
            UserRole::DOCTOR      => 'Doctor',
            UserRole::ASSISTANT   => 'Assistant',
        };
    }

    public function isClinicalStaff(): bool
    {
        return match($this) {
            UserRole::SUPER_ADMIN => false,
            UserRole::DOCTOR,
            UserRole::ASSISTANT   => true,
        };
    }
}
\end{lstlisting}

\begin{notebox}[Enum Casting]
  Because \texttt{UserRole} is declared in the \texttt{casts()} method
  of the \phpclass{User} model, Eloquent automatically deserialises the
  raw string value stored in the database
  (\texttt{"super\_admin"}, \texttt{"doctor"}, \texttt{"assistant"})
  into a \phpclass{UserRole} enum instance when a \phpclass{User} model
  is hydrated. Comparisons throughout the application therefore use the
  type-safe \texttt{UserRole::DOCTOR} form rather than the raw string
  \texttt{"doctor"}, eliminating an entire class of typo-induced bugs.
\end{notebox}

\subsection{AuthService: Login Logic}

The \phpclass{AuthService} class encapsulates all authentication
business logic. The controller delegates to this service entirely and
contains no credential-checking code of its own.

\begin{lstlisting}[style=php, caption={\texttt{app/Services/AuthService.php} ---
  Credential verification, active account check, and token issuance.}]
<?php

namespace App\Services;

use App\Exceptions\AccountDeactivatedException;
use App\Models\User;
use Illuminate\Support\Facades\Auth;

class AuthService
{
    public function login(string $email, string $password): array
    {
        if (! Auth::attempt(['email' => $email, 'password' => $password])) {
            throw new \Illuminate\Validation\ValidationException(
                validator([], []),
                response()->json([
                    'success' => false,
                    'message' => 'The provided credentials are incorrect.',
                ], 401)
            );
        }

        /** @var User $user */
        $user = Auth::user();

        if (! $user->is_active) {
            Auth::logout();
            throw new AccountDeactivatedException();
        }

        $token = $user->createToken('clinic-management-api')->accessToken;

        return [
            'user'  => $user,
            'token' => $token,
        ];
    }

    public function logout(User $user): void
    {
        $user->token()->revoke();
    }

    public function me(User $user): User
    {
        return $user->load('clinic');
    }
}
\end{lstlisting}

\begin{warningbox}[AccountDeactivatedException]
  After a successful credential match, the service performs a second
  check against the \texttt{is\_active} flag before issuing a token.
  This ordering is intentional: credentials are verified first so that
  the error message does not reveal whether an email address exists in
  the system. Only after credentials are confirmed valid does the system
  disclose that the account is deactivated. This prevents user
  enumeration via the login endpoint.
\end{warningbox}

\subsection{AuthController: HTTP Layer}

The \phpclass{AuthController} contains no business logic. It receives
the validated form request, delegates to the service, and serialises the
response.

\begin{lstlisting}[style=php, caption={\texttt{app/Http/Controllers/AuthController.php}
  --- Thin controller delegating entirely to \phpclass{AuthService}.}]
<?php

namespace App\Http\Controllers;

use App\Http\Requests\Auth\LoginRequest;
use App\Http\Resources\UserResource;
use App\Services\AuthService;
use App\Traits\ApiResponse;
use Illuminate\Http\JsonResponse;
use Illuminate\Http\Request;

class AuthController extends Controller
{
    use ApiResponse;

    public function __construct(
        private readonly AuthService $authService,
    ) {}

    public function login(LoginRequest $request): JsonResponse
    {
        $result = $this->authService->login(
            email:    $request->validated('email'),
            password: $request->validated('password'),
        );

        return $this->successResponse(
            data: [
                'user'         => new UserResource($result['user']),
                'access_token' => $result['token'],
                'token_type'   => 'Bearer',
            ],
            message: 'Login successful.',
            status: 200,
        );
    }

    public function logout(Request $request): JsonResponse
    {
        $this->authService->logout($request->user());

        return $this->successResponse(
            data:    [],
            message: 'Logged out successfully.',
        );
    }

    public function me(Request $request): JsonResponse
    {
        $user = $this->authService->me($request->user());

        return $this->successResponse(
            data:    new UserResource($user),
            message: 'Authenticated user retrieved.',
        );
    }
}
\end{lstlisting}

\subsection{Login Request Validation}

The \phpclass{LoginRequest} form request enforces structure before the
controller is even reached.

\begin{lstlisting}[style=php, caption={\texttt{app/Http/Requests/Auth/LoginRequest.php}.}]
<?php

namespace App\Http\Requests\Auth;

use Illuminate\Foundation\Http\FormRequest;

class LoginRequest extends FormRequest
{
    public function authorize(): bool
    {
        return true;
    }

    public function rules(): array
    {
        return [
            'email'    => ['required', 'string', 'email'],
            'password' => ['required', 'string'],
        ];
    }
}
\end{lstlisting}

\subsection{Token Issuance and the oauth\_access\_tokens Table}

When \texttt{\$user->createToken('clinic-management-api')} is called,
Passport performs the following sequence:

\begin{enumerate}
  \item Generates a cryptographically random 80-character token string.
  \item Splits the string: the first 40 characters are the
    \textbf{token ID}; the second 40 characters are the
    \textbf{token secret}.
  \item Hashes the secret with SHA-256 and persists a record to
    \texttt{oauth\_access\_tokens} containing the token ID, the hashed
    secret, the user ID, the client ID, and expiry metadata.
  \item Returns the raw (unhashed) 80-character string to the
    application, which is returned to the client only once. It is never
    stored in plaintext anywhere in the system.
\end{enumerate}

On every subsequent request bearing this token, Passport's
\texttt{auth:api} middleware repeats the split, hashes the secret
portion, and queries \texttt{oauth\_access\_tokens} by token ID for a
record whose hashed secret matches and whose \texttt{revoked} flag is
\texttt{0}. Only when both conditions hold is the request considered
authenticated.

\begin{table}[h!]
  \centering
  \caption{Key columns in the \texttt{oauth\_access\_tokens} table
    managed by Laravel Passport.}
  \label{tab:oauth-tokens-table}
  \renewcommand{\arraystretch}{1.4}
  \begin{tabularx}{\linewidth}{@{}L{3.5cm} L{2.8cm} X@{}}
    \toprule
    \rowcolor{PrimaryBlue}
    \tblhdr{Column} & \tblhdr{Type} & \tblhdr{Purpose} \\
    \midrule
    \texttt{id}          & \texttt{VARCHAR(100)} & The first 40 chars of the issued
                                                   token; used as the lookup key. \\
    \rowalt
    \texttt{user\_id}    & \texttt{CHAR(36)}     & UUID of the owning \phpclass{User}
                                                   record. \\
    \texttt{client\_id}  & \texttt{CHAR(36)}     & The Passport client that issued
                                                   this token. \\
    \rowalt
    \texttt{name}        & \texttt{VARCHAR(255)} & Human-readable label passed to
                                                   \texttt{createToken()}. \\
    \texttt{scopes}      & \texttt{TEXT}         & JSON array of OAuth2 scopes
                                                   (empty array for PATs in this API). \\
    \rowalt
    \texttt{revoked}     & \texttt{TINYINT(1)}   & Set to \texttt{1} on logout;
                                                   causes all future requests to
                                                   be rejected. \\
    \texttt{expires\_at} & \texttt{DATETIME}     & Token expiry timestamp. Configurable
                                                   via \texttt{Passport::tokensExpireIn()}. \\
    \bottomrule
  \end{tabularx}
\end{table}

\subsection{Token Revocation on Logout}

Token revocation is deliberately simple. When the authenticated user
calls \route{POST /api/\apiversion{}/auth/logout}, the service calls
\texttt{\$user->token()->revoke()}, which executes:

\begin{lstlisting}[style=php, caption={Passport's internal revocation update.}]
DB::table('oauth_access_tokens')
    ->where('id', $this->id)
    ->update(['revoked' => true]);
\end{lstlisting}

From this point forward, the \texttt{auth:api} middleware will find
\texttt{revoked = 1} for this token and return a \stunauth{} response.
The client is expected to discard the token from its local storage.

%% ─────────────────────────────────────────────────────────────────────────
\section{Canonical JSON Request and Response Examples}
\label{sec:auth-examples}
%% ─────────────────────────────────────────────────────────────────────────

\subsection{Login}

\begin{lstlisting}[style=json, caption={Request: \texttt{POST
  /api/\apiversion{}/auth/login}.}]
{
  "email": "doctor@clinic-a.example.com",
  "password": "secret1234"
}
\end{lstlisting}

\begin{lstlisting}[style=json, caption={Response: HTTP 200 OK on
  successful authentication.}]
{
  "success": true,
  "message": "Login successful.",
  "data": {
    "user": {
      "id": "9d2b1f4e-3a7c-4e8d-b1f2-3a7c4e8db1f2",
      "name": "Dr. Sarah Khalil",
      "email": "doctor@clinic-a.example.com",
      "role": "doctor",
      "is_active": true,
      "clinic_id": "1a2b3c4d-5e6f-7a8b-9c0d-1e2f3a4b5c6d"
    },
    "access_token": "eyJ0eXAiOiJKV1Qi...<truncated>",
    "token_type": "Bearer"
  }
}
\end{lstlisting}

\subsection{Using the Token}

Every subsequent request must present the token in the
\texttt{Authorization} header:

\begin{lstlisting}[style=bash, caption={Bearer token usage in a curl request.}]
curl -X GET https://api.example.com/api/v1/doctor/appointments \
  -H "Accept: application/json" \
  -H "Authorization: Bearer eyJ0eXAiOiJKV1Qi...<token>"
\end{lstlisting}

\subsection{Unauthenticated Response}

A request that omits the \texttt{Authorization} header, presents a
malformed token, or presents a revoked token receives:

\begin{lstlisting}[style=json, caption={Response: HTTP 401 Unauthorized.}]
{
  "success": false,
  "message": "Unauthenticated."
}
\end{lstlisting}

%% ─────────────────────────────────────────────────────────────────────────
\section{The Role-Based Access Control System}
\label{sec:rbac}
%% ─────────────────────────────────────────────────────────────────────────

\subsection{Role Hierarchy Definition}

The application defines three discrete, non-overlapping roles. There is
no role inheritance; a user can hold exactly one role at any given time,
and the permission surface of that role is strictly bounded.

\begin{table}[h!]
  \centering
  \caption{Role hierarchy, scope, and permission surface summary.}
  \label{tab:role-hierarchy}
  \renewcommand{\arraystretch}{1.5}
  \begin{tabularx}{\linewidth}{@{}L{3cm} L{2.8cm} X@{}}
    \toprule
    \rowcolor{PrimaryBlue}
    \tblhdr{Role} & \tblhdr{Enum Value} & \tblhdr{Scope and Summary} \\
    \midrule
    Super Administrator
      & \texttt{super\_admin}
      & Global scope. Manages clinics and all user accounts.
        No access to clinical data (appointments, patients,
        prescriptions). \\
    \rowalt
    Doctor
      & \texttt{doctor}
      & Clinic-scoped. Manages own schedule, views own
        appointments, advances appointment status, authors
        prescriptions. Read-only access to patients. \\
    Assistant
      & \texttt{assistant}
      & Clinic-scoped. Registers and updates patients, books
        and reschedules appointments, cancels appointments,
        queries available slots. No access to prescriptions. \\
    \bottomrule
  \end{tabularx}
\end{table}

\begin{importantbox}[Non-Overlapping Permission Surfaces]
  The role surfaces are deliberately non-overlapping by design. A Super
  Administrator who authenticates cannot access
  \route{/api/\apiversion{}/doctor/} or
  \route{/api/\apiversion{}/assistant/} routes. A Doctor who attempts
  to access \route{/api/\apiversion{}/admin/} routes will receive a
  \stforbid{} response. This strict separation ensures that a
  compromised Super Administrator account cannot be used to forge
  prescriptions or alter patient records, limiting the blast radius of
  any single credential compromise.
\end{importantbox}

\subsection{The RoleMiddleware Implementation}

The \phpclass{RoleMiddleware} is the primary enforcement point for
role-based access control. It is applied at the route-group level in
\texttt{routes/api.php}, meaning it executes after \texttt{auth:api}
has already confirmed the token is valid and populated
\texttt{request()->user()}.

\begin{lstlisting}[style=php, caption={\texttt{app/Http/Middleware/RoleMiddleware.php}
  --- Strict role gate using the UserRole enum.}]
<?php

namespace App\Http\Middleware;

use App\Enums\UserRole;
use Closure;
use Illuminate\Http\JsonResponse;
use Illuminate\Http\Request;

class RoleMiddleware
{
    public function handle(Request $request, Closure $next, string ...$roles): JsonResponse|mixed
    {
        $user = $request->user();

        $allowedRoles = array_map(
            fn(string $role) => UserRole::from($role),
            $roles,
        );

        if (! in_array($user->role, $allowedRoles, strict: true)) {
            return response()->json([
                'success' => false,
                'message' => 'You do not have permission to perform this action.',
            ], 403);
        }

        return $next($request);
    }
}
\end{lstlisting}

\begin{notebox}[Strict Enum Comparison]
  The \texttt{in\_array} call uses \texttt{strict: true}. Because
  \texttt{\$user->role} is a \phpclass{UserRole} enum instance (due to
  Eloquent casting) and \texttt{\$allowedRoles} is an array of
  \phpclass{UserRole} enum instances, the comparison uses PHP's
  built-in enum identity semantics. This eliminates any possibility of
  a type-coercion bypass where a string \texttt{"doctor"} might
  accidentally match a different type.
\end{notebox}

\subsection{Route Group Structure and Middleware Stacking}

The three role surfaces are implemented as three separate route groups
in \texttt{routes/api.php}. Every route within a group inherits the
entire middleware stack of that group.

\begin{lstlisting}[style=php, caption={\texttt{routes/api.php} ---
  Route group definitions with stacked middleware.}]
<?php

use App\Http\Controllers\AuthController;
use App\Http\Controllers\Admin\ClinicController;
use App\Http\Controllers\Admin\UserController;
use App\Http\Controllers\Admin\DashboardController;
use App\Http\Controllers\Doctor\AppointmentController as DoctorApptController;
use App\Http\Controllers\Doctor\DoctorScheduleController;
use App\Http\Controllers\Doctor\PatientController as DoctorPatientController;
use App\Http\Controllers\Doctor\PrescriptionController;
use App\Http\Controllers\Assistant\AppointmentController as AssistantApptController;
use App\Http\Controllers\Assistant\PatientController as AssistantPatientController;
use App\Http\Controllers\Assistant\AvailableSlotController;
use Illuminate\Support\Facades\Route;

/*
|--------------------------------------------------------------------------
| Public Routes (no authentication required)
|--------------------------------------------------------------------------
*/
Route::prefix('v1')->group(function () {

    Route::prefix('auth')->group(function () {
        Route::post('login',  [AuthController::class, 'login']);
    });

    /*
    |----------------------------------------------------------------------
    | Authenticated Routes
    |----------------------------------------------------------------------
    */
    Route::middleware('auth:api')->group(function () {

        Route::prefix('auth')->group(function () {
            Route::post('logout', [AuthController::class, 'logout']);
            Route::get('me',      [AuthController::class, 'me']);
        });

        /*
        |------------------------------------------------------------------
        | Super Administrator Routes
        |------------------------------------------------------------------
        */
        Route::prefix('admin')
            ->middleware('role:super_admin')
            ->group(function () {

            Route::get('dashboard', [DashboardController::class, 'index']);

            Route::apiResource('clinics', ClinicController::class);

            Route::apiResource('users', UserController::class);
            Route::patch('users/{user}/role',
                [UserController::class, 'updateRole']);
            Route::patch('users/{user}/status',
                [UserController::class, 'updateStatus']);
        });

        /*
        |------------------------------------------------------------------
        | Doctor Routes
        |------------------------------------------------------------------
        */
        Route::prefix('doctor')
            ->middleware(['role:doctor', 'clinic.scope'])
            ->group(function () {

            Route::apiResource('schedules', DoctorScheduleController::class);

            Route::get('appointments',
                [DoctorApptController::class, 'index']);
            Route::get('appointments/{appointment}',
                [DoctorApptController::class, 'show']);
            Route::patch('appointments/{appointment}/status',
                [DoctorApptController::class, 'updateStatus']);

            Route::get('patients',
                [DoctorPatientController::class, 'index']);
            Route::get('patients/{patient}',
                [DoctorPatientController::class, 'show']);

            Route::apiResource('prescriptions', PrescriptionController::class)
                ->except(['index', 'destroy']);
        });

        /*
        |------------------------------------------------------------------
        | Assistant Routes
        |------------------------------------------------------------------
        */
        Route::prefix('assistant')
            ->middleware(['role:assistant', 'clinic.scope'])
            ->group(function () {

            Route::get('available-slots',
                [AvailableSlotController::class, 'index']);

            Route::apiResource('patients',
                AssistantPatientController::class)->except(['destroy']);

            Route::get('appointments',
                [AssistantApptController::class, 'index']);
            Route::get('appointments/{appointment}',
                [AssistantApptController::class, 'show']);
            Route::post('appointments',
                [AssistantApptController::class, 'store']);
            Route::put('appointments/{appointment}',
                [AssistantApptController::class, 'update']);
            Route::delete('appointments/{appointment}',
                [AssistantApptController::class, 'destroy']);
        });
    });
});
\end{lstlisting}

%% ─────────────────────────────────────────────────────────────────────────
\section{Clinic Scope Middleware}
\label{sec:clinic-scope-middleware}
%% ─────────────────────────────────────────────────────────────────────────

The \phpclass{ClinicScopeMiddleware} runs on every Doctor and Assistant
route. Its sole responsibility is to inject the authenticated user's
\texttt{clinic\_id} into the request data bag, overwriting any value
the client may have provided.

\begin{lstlisting}[style=php, caption={\texttt{app/Http/Middleware/ClinicScopeMiddleware.php}
  --- Overwrite client-supplied clinic\_id with the authenticated
  user's authoritative value.}]
<?php

namespace App\Http\Middleware;

use Closure;
use Illuminate\Http\Request;

class ClinicScopeMiddleware
{
    public function handle(Request $request, Closure $next): mixed
    {
        $user = $request->user();

        // Overwrite any client-supplied clinic_id unconditionally.
        // This prevents cross-clinic data forgery via request body manipulation.
        $request->merge(['clinic_id' => $user->clinic_id]);

        return $next($request);
    }
}
\end{lstlisting}

%% ─────────────────────────────────────────────────────────────────────────
\section{Comprehensive Endpoint-to-Role Mapping}
\label{sec:endpoint-role-map}
%% ─────────────────────────────────────────────────────────────────────────

The following table catalogues every endpoint in the API and specifies
the exact authentication and authorisation requirements for each. The
column headings use the following abbreviations: \textbf{SA} = Super
Administrator (\texttt{super\_admin}), \textbf{DR} = Doctor
(\texttt{doctor}), \textbf{AS} = Assistant (\texttt{assistant}),
\textbf{PUB} = Public (no authentication required).

A tick (\checkmark) indicates the role is permitted. A dash (---) in
the Role column indicates no authentication is required. Clinic scope
enforcement (CS) is indicated where the \phpclass{ClinicScopeMiddleware}
is active in addition to the role gate.

\renewcommand{\arraystretch}{1.35}

\begin{longtable}{@{}L{1.1cm} L{5.4cm} L{2.0cm} C{1.0cm} C{1.0cm} C{1.0cm} C{0.8cm}@{}}
  \toprule
  \rowcolor{PrimaryBlue}
  \tblhdr{Method} &
  \tblhdr{Endpoint} &
  \tblhdr{Action} &
  \tblhdr{SA} &
  \tblhdr{DR} &
  \tblhdr{AS} &
  \tblhdr{PUB} \\
  \midrule
  \endfirsthead

  \toprule
  \rowcolor{PrimaryBlue}
  \tblhdr{Method} &
  \tblhdr{Endpoint (cont.)} &
  \tblhdr{Action} &
  \tblhdr{SA} &
  \tblhdr{DR} &
  \tblhdr{AS} &
  \tblhdr{PUB} \\
  \midrule
  \endhead

  \bottomrule
  \endfoot

  \multicolumn{7}{@{}l}{\small\textit{--- Authentication ---}} \\[2pt]
  \mpost    & \route{/api/v1/auth/login}         & Obtain bearer token        & --- & --- & --- & \checkmark \\
  \rowalt
  \mpost    & \route{/api/v1/auth/logout}         & Revoke current token       & \checkmark & \checkmark & \checkmark & \\
  \mget     & \route{/api/v1/auth/me}             & Retrieve own profile       & \checkmark & \checkmark & \checkmark & \\
  \rowalt
  \\[-6pt]
  \multicolumn{7}{@{}l}{\small\textit{--- Super Administrator: Dashboard ---}} \\[2pt]
  \mget     & \route{/api/v1/admin/dashboard}     & System KPI aggregate       & \checkmark & & & \\
  \rowalt
  \\[-6pt]
  \multicolumn{7}{@{}l}{\small\textit{--- Super Administrator: Clinic Management ---}} \\[2pt]
  \mget     & \route{/api/v1/admin/clinics}               & List all clinics           & \checkmark & & & \\
  \rowalt
  \mpost    & \route{/api/v1/admin/clinics}               & Create clinic              & \checkmark & & & \\
  \mget     & \route{/api/v1/admin/clinics/\{id\}}        & Retrieve clinic            & \checkmark & & & \\
  \rowalt
  \mput     & \route{/api/v1/admin/clinics/\{id\}}        & Update clinic              & \checkmark & & & \\
  \mdelete  & \route{/api/v1/admin/clinics/\{id\}}        & Delete clinic              & \checkmark & & & \\
  \rowalt
  \\[-6pt]
  \multicolumn{7}{@{}l}{\small\textit{--- Super Administrator: User Management ---}} \\[2pt]
  \mget     & \route{/api/v1/admin/users}                 & List all users             & \checkmark & & & \\
  \rowalt
  \mpost    & \route{/api/v1/admin/users}                 & Create user                & \checkmark & & & \\
  \mget     & \route{/api/v1/admin/users/\{id\}}          & Retrieve user              & \checkmark & & & \\
  \rowalt
  \mput     & \route{/api/v1/admin/users/\{id\}}          & Update user                & \checkmark & & & \\
  \mdelete  & \route{/api/v1/admin/users/\{id\}}          & Delete user                & \checkmark & & & \\
  \rowalt
  \mpatch   & \route{/api/v1/admin/users/\{id\}/role}     & Change user role           & \checkmark & & & \\
  \mpatch   & \route{/api/v1/admin/users/\{id\}/status}   & Activate / deactivate user & \checkmark & & & \\
  \rowalt
  \\[-6pt]
  \multicolumn{7}{@{}l}{\small\textit{--- Doctor: Schedule Management (CS enforced) ---}} \\[2pt]
  \mget     & \route{/api/v1/doctor/schedules}            & List own schedules         & & \checkmark & & \\
  \rowalt
  \mpost    & \route{/api/v1/doctor/schedules}            & Create schedule entry      & & \checkmark & & \\
  \mget     & \route{/api/v1/doctor/schedules/\{id\}}     & Retrieve schedule entry    & & \checkmark & & \\
  \rowalt
  \mput     & \route{/api/v1/doctor/schedules/\{id\}}     & Update schedule entry      & & \checkmark & & \\
  \mdelete  & \route{/api/v1/doctor/schedules/\{id\}}     & Delete schedule entry      & & \checkmark & & \\
  \rowalt
  \\[-6pt]
  \multicolumn{7}{@{}l}{\small\textit{--- Doctor: Appointment Management (CS enforced) ---}} \\[2pt]
  \mget     & \route{/api/v1/doctor/appointments}                    & List own appointments      & & \checkmark & & \\
  \rowalt
  \mget     & \route{/api/v1/doctor/appointments/\{id\}}             & Retrieve appointment       & & \checkmark & & \\
  \mpatch   & \route{/api/v1/doctor/appointments/\{id\}/status}      & Confirm or complete appt.  & & \checkmark & & \\
  \rowalt
  \\[-6pt]
  \multicolumn{7}{@{}l}{\small\textit{--- Doctor: Patient Read Access (CS enforced) ---}} \\[2pt]
  \mget     & \route{/api/v1/doctor/patients}             & List clinic patients       & & \checkmark & & \\
  \rowalt
  \mget     & \route{/api/v1/doctor/patients/\{id\}}      & Retrieve patient record    & & \checkmark & & \\
  \\[-6pt]
  \rowalt
  \multicolumn{7}{@{}l}{\small\textit{--- Doctor: Prescription Authoring (CS enforced) ---}} \\[2pt]
  \mpost    & \route{/api/v1/doctor/prescriptions}        & Create prescription        & & \checkmark & & \\
  \rowalt
  \mget     & \route{/api/v1/doctor/prescriptions/\{id\}} & Retrieve prescription      & & \checkmark & & \\
  \mput     & \route{/api/v1/doctor/prescriptions/\{id\}} & Update prescription        & & \checkmark & & \\
  \rowalt
  \\[-6pt]
  \multicolumn{7}{@{}l}{\small\textit{--- Assistant: Slot Availability (CS enforced) ---}} \\[2pt]
  \mget     & \route{/api/v1/assistant/available-slots}              & Query available slots      & & & \checkmark & \\
  \rowalt
  \\[-6pt]
  \multicolumn{7}{@{}l}{\small\textit{--- Assistant: Patient Management (CS enforced) ---}} \\[2pt]
  \mget     & \route{/api/v1/assistant/patients}          & List clinic patients       & & & \checkmark & \\
  \rowalt
  \mpost    & \route{/api/v1/assistant/patients}          & Register new patient       & & & \checkmark & \\
  \mget     & \route{/api/v1/assistant/patients/\{id\}}   & Retrieve patient record    & & & \checkmark & \\
  \rowalt
  \mput     & \route{/api/v1/assistant/patients/\{id\}}   & Update patient record      & & & \checkmark & \\
  \\[-6pt]
  \rowalt
  \multicolumn{7}{@{}l}{\small\textit{--- Assistant: Appointment Management (CS enforced) ---}} \\[2pt]
  \mget     & \route{/api/v1/assistant/appointments}                 & List clinic appointments   & & & \checkmark & \\
  \rowalt
  \mget     & \route{/api/v1/assistant/appointments/\{id\}}          & Retrieve appointment       & & & \checkmark & \\
  \mpost    & \route{/api/v1/assistant/appointments}                  & Book appointment           & & & \checkmark & \\
  \rowalt
  \mput     & \route{/api/v1/assistant/appointments/\{id\}}           & Reschedule appointment     & & & \checkmark & \\
  \mdelete  & \route{/api/v1/assistant/appointments/\{id\}}           & Cancel appointment         & & & \checkmark & \\
\end{longtable}

\begin{notebox}[Clinic Scope (CS) Enforcement]
  Every row in the Doctor and Assistant sections carries implicit clinic
  scope enforcement via \phpclass{ClinicScopeMiddleware}. Even if an
  endpoint is listed as accessible to a Doctor, a Doctor from Clinic~A
  cannot access resources belonging to Clinic~B. Ownership is validated
  at the controller layer for resource-by-UUID endpoints and at the
  query layer (via an automatic \texttt{where('clinic\_id', ...)} scope)
  for list endpoints.
\end{notebox}

%% ─────────────────────────────────────────────────────────────────────────
\section{Privilege Escalation: Threat Model and Defences}
\label{sec:privilege-escalation}
%% ─────────────────────────────────────────────────────────────────────────

Privilege escalation is the class of attacks in which an authenticated
principal gains access to resources or capabilities beyond those
explicitly granted to their role. The two canonical forms are
\textbf{vertical privilege escalation} (gaining a higher-privilege
role's capabilities) and \textbf{horizontal privilege escalation}
(accessing another principal's data within the same role).

\subsection{Vertical Privilege Escalation}

Vertical privilege escalation occurs when a lower-privileged user
successfully invokes an endpoint that requires a higher-privileged role,
or when a user is able to elevate their own role.

\subsubsection{Attack Vector 1: Role Forgery via Request Body}

An attacker who authenticates as an Assistant might craft a request body
that includes \texttt{"role": "super\_admin"}, attempting to trick the
application into treating the request as if it originated from a higher
role.

\textbf{Defence.} The \phpclass{RoleMiddleware} reads the role
exclusively from \texttt{\$request->user()->role}, which is the
\texttt{UserRole} enum hydrated from the database record belonging to
the authenticated token owner. The request body is never consulted for
role determination. A \texttt{"role"} field in the request body is
simply ignored or causes a validation error if the endpoint does not
accept it.

\subsubsection{Attack Vector 2: Self-Role Elevation via the Admin API}

A Super Administrator account could theoretically be used to elevate a
lower-privileged account's role if there were no guards on the role
change endpoint. More critically, a Super Administrator could attempt to
demote their own account to a Doctor and re-promote it to bypass audit
logs, or could be tricked into demoting their own account, locking
themselves (and potentially the organisation) out of administrative
access.

\textbf{Defence.} The \phpclass{UserController}'s \texttt{updateRole}
action contains an explicit self-demotion guard:

\begin{lstlisting}[style=php, caption={Self-demotion guard in
  \texttt{UserController::updateRole()}.}]
public function updateRole(UpdateUserRoleRequest $request, User $user): JsonResponse
{
    if ($user->id === $request->user()->id) {
        throw new SelfDemotionException();
    }

    $updatedUser = $this->userService->updateRole(
        user: $user,
        role: UserRole::from($request->validated('role')),
    );

    return $this->successResponse(
        data:    new UserResource($updatedUser),
        message: 'User role updated successfully.',
    );
}
\end{lstlisting}

The \phpclass{SelfDemotionException} renders as a \stforbid{} response
with the message \textit{``You cannot change your own role.''} This
prevents accidental or malicious self-demotion while still allowing a
Super Administrator to manage all other accounts freely.

\subsubsection{Attack Vector 3: Route Traversal Across Role Prefixes}

An authenticated Doctor might construct a request to a Super
Administrator endpoint such as \route{POST /api/\apiversion{}/admin/users},
hoping that the authentication middleware alone would admit the request.

\textbf{Defence.} The \phpclass{RoleMiddleware} is applied at the route
group level, meaning it executes for every route within the group
regardless of the HTTP method or path specifics. A Doctor's token is
valid for the \texttt{auth:api} middleware but fails the
\texttt{role:super\_admin} check and receives a \stforbid{} response.
The role check happens before the controller is invoked, so the attack
is terminated in the middleware pipeline.

\begin{table}[h!]
  \centering
  \caption{Vertical privilege escalation attack vectors and their mitigations.}
  \label{tab:vert-escalation}
  \renewcommand{\arraystretch}{1.5}
  \begin{tabularx}{\linewidth}{@{}L{4.5cm} X@{}}
    \toprule
    \rowcolor{PrimaryBlue}
    \tblhdr{Attack Vector} & \tblhdr{Mitigation} \\
    \midrule
    Role forgery via request body
      & Role read exclusively from \texttt{\$request->user()->role};
        request body \texttt{role} field is ignored. \\
    \rowalt
    Self-demotion / self-role-change
      & Explicit \texttt{\$user->id === \$request->user()->id} guard in
        \texttt{updateRole()}; throws \phpclass{SelfDemotionException}
        (\stforbid{}). \\
    Route traversal across role prefixes
      & \phpclass{RoleMiddleware} applied at route-group level; executes
        before any controller code is reached. \\
    \rowalt
    Credential theft with deactivated account
      & \texttt{is\_active} checked after credential verification in
        \phpclass{AuthService}; \phpclass{AccountDeactivatedException}
        thrown (\stforbid{}). \\
    Token reuse after logout
      & Passport's server-side revocation sets \texttt{revoked = 1} in
        \texttt{oauth\_access\_tokens}; subsequent requests with the
        same token receive \stunauth{}. \\
    \bottomrule
  \end{tabularx}
\end{table}

\subsection{Horizontal Privilege Escalation}

Horizontal privilege escalation occurs when an authenticated user
accesses, modifies, or deletes data belonging to another principal
\emph{within the same role level}. In the context of this API, the two
most critical horizontal escalation surfaces are cross-clinic data access
and cross-doctor prescription access.

\subsubsection{Attack Vector 4: Cross-Clinic Resource Access by UUID Guessing}

All primary keys in the system are UUID v4 values. A Doctor at Clinic~A
might attempt to enumerate patient records at Clinic~B by guessing or
obtaining UUIDs (e.g.\ from a shared infrastructure log) and supplying
them to the \route{GET /api/\apiversion{}/doctor/patients/\{id\}} endpoint.

\textbf{Defence.} Laravel's route model binding resolves the UUID to an
Eloquent model instance. The controller then compares the resolved model's
\texttt{clinic\_id} against the authenticated user's \texttt{clinic\_id}
(which was injected by \phpclass{ClinicScopeMiddleware}):

\begin{lstlisting}[style=php, caption={Cross-clinic ownership check in a
  Doctor route controller.}]
public function show(Request $request, Patient $patient): JsonResponse
{
    if ($patient->clinic_id !== $request->input('clinic_id')) {
        throw new ClinicScopeViolationException();
    }

    return $this->successResponse(
        data:    new PatientResource($patient),
        message: 'Patient retrieved successfully.',
    );
}
\end{lstlisting}

The \phpclass{ClinicScopeViolationException} renders as a \stforbid{}
response. Critically, the response does not differ from a genuine
\stforbid{} triggered by a role mismatch, so the attacker cannot
distinguish between ``this resource does not exist'' and ``this resource
exists but you cannot access it.'' This prevents oracle-style
enumeration.

\subsubsection{Attack Vector 5: Cross-Doctor Prescription Access}

A prescription is authored by a specific Doctor and is associated with
a specific appointment. A different Doctor at the same clinic should not
be able to view or modify another Doctor's prescription, even though
they share the same role and the same clinic scope.

\textbf{Defence.} The \phpclass{PrescriptionController} implements an
additional ownership check beyond the clinic scope gate:

\begin{lstlisting}[style=php, caption={Cross-doctor ownership check in
  \texttt{PrescriptionController}.}]
public function show(Request $request, Prescription $prescription): JsonResponse
{
    // Layer 1: clinic scope (enforced by ClinicScopeMiddleware + check)
    if ($prescription->appointment->clinic_id !== $request->input('clinic_id')) {
        throw new ClinicScopeViolationException();
    }

    // Layer 2: doctor ownership (only the authoring doctor may view)
    if ($prescription->appointment->doctor_id !== $request->user()->id) {
        throw new ClinicScopeViolationException();
    }

    return $this->successResponse(
        data:    new PrescriptionResource($prescription->load('items')),
        message: 'Prescription retrieved successfully.',
    );
}
\end{lstlisting}

\subsubsection{Attack Vector 6: Cross-Clinic Data Injection via Request
  Body Forgery}

An Assistant at Clinic~A might craft a \route{POST} request to
\route{/api/\apiversion{}/assistant/patients} with
\texttt{"clinic\_id": "<uuid-of-clinic-b>"} in the JSON body, attempting
to create a patient record in a different clinic's namespace.

\textbf{Defence.} The \phpclass{ClinicScopeMiddleware} calls
\texttt{\$request->merge(['clinic\_id' => \$user->clinic\_id])} before
the request reaches the form request validator or the controller. The
\texttt{merge()} call overwrites any existing \texttt{clinic\_id} key
in the request data bag with the authoritative value from the
authenticated user's database record. The client-supplied value is
silently discarded.

\subsubsection{Attack Vector 7: Appointment Double-Booking via Race Condition}

Two Assistants at the same clinic might simultaneously submit booking
requests for the same doctor, date, and time slot. Without a
synchronisation mechanism, both requests could pass the slot availability
check (both read an empty slot) and then both write a record, resulting
in two patients occupying the same appointment.

\textbf{Defence.} The \phpclass{AppointmentService} acquires a
distributed cache lock before performing the slot check and write:

\begin{lstlisting}[style=php, caption={Distributed lock in
  \texttt{AppointmentService::book()} to prevent double-booking.}]
$lockKey = "appointment_slot:{$doctorId}:{$date}:{$startTime}";

$lock = Cache::lock($lockKey, seconds: 10);

if (! $lock->get()) {
    throw new AppointmentConflictException();
}

try {
    $this->validateSlotAvailability($doctorId, $date, $startTime);
    $appointment = $this->appointmentRepository->create($data);
    $lock->release();
    return $appointment;
} catch (\Throwable $e) {
    $lock->release();
    throw $e;
}
\end{lstlisting}

As a second line of defence, the database enforces a unique composite
index on \texttt{(doctor\_id, appointment\_date, start\_time)}, so even
if the lock is somehow bypassed, the database write will fail with a
unique constraint violation, which the repository layer converts to an
\phpclass{AppointmentConflictException} (\stconflict{}).

\begin{table}[h!]
  \centering
  \caption{Horizontal privilege escalation attack vectors and their mitigations.}
  \label{tab:horiz-escalation}
  \renewcommand{\arraystretch}{1.5}
  \begin{tabularx}{\linewidth}{@{}L{4.5cm} X@{}}
    \toprule
    \rowcolor{PrimaryBlue}
    \tblhdr{Attack Vector} & \tblhdr{Mitigation} \\
    \midrule
    Cross-clinic resource access by UUID
      & Explicit \texttt{clinic\_id} ownership check in every
        resource-by-UUID controller action; throws
        \phpclass{ClinicScopeViolationException} (\stforbid{}).
        Response is indistinguishable from a role mismatch to prevent
        oracle enumeration. \\
    \rowalt
    Cross-doctor prescription access
      & Secondary \texttt{doctor\_id} ownership check in
        \phpclass{PrescriptionController} beyond the clinic scope
        gate; same exception and response. \\
    Cross-clinic data injection via body
      & \phpclass{ClinicScopeMiddleware} overwrites client-supplied
        \texttt{clinic\_id} with the authenticated user's authoritative
        value before the controller is reached. \\
    \rowalt
    Appointment double-booking race condition
      & Distributed \texttt{Cache::lock()} around the slot check and
        write in \phpclass{AppointmentService}; unique composite index
        on \texttt{(doctor\_id, appointment\_date, start\_time)} as
        final database-level backstop. \\
    Patient registration in another clinic
      & Same as cross-clinic body injection above; \texttt{clinic\_id}
        is always sourced from the authenticated user, never from the
        client payload. \\
    \bottomrule
  \end{tabularx}
\end{table}

%% ─────────────────────────────────────────────────────────────────────────
\section{The AccountDeactivatedException Flow}
\label{sec:deactivated-account}
%% ─────────────────────────────────────────────────────────────────────────

User accounts carry an \texttt{is\_active} boolean flag that can be set
to \texttt{false} by a Super Administrator via the
\route{PATCH /api/\apiversion{}/admin/users/\{id\}/status} endpoint.
Deactivation is a soft mechanism: the user record is retained in the
database and all historical data (appointments, prescriptions) remains
intact and attributable to the account. Only the ability to obtain a new
token is blocked.

It is important to note that deactivation alone does not revoke tokens
that were issued \emph{before} the account was deactivated. Revoking
pre-existing tokens requires a separate administrative action: the Super
Administrator must also call the token revocation endpoint (or, if not
yet implemented, the \texttt{oauth\_access\_tokens} table must be
updated directly to set \texttt{revoked = 1} for all tokens belonging
to the user).

\begin{warningbox}[Token Revocation on Deactivation]
  The current implementation blocks new token issuance for deactivated
  accounts but does not automatically revoke existing tokens. If
  immediate session termination is required (e.g.\ a staff member is
  dismissed), the Super Administrator must both deactivate the account
  and explicitly trigger token revocation. A future enhancement should
  atomically revoke all active tokens when \texttt{is\_active} is set
  to \texttt{false}.
\end{warningbox}

\begin{lstlisting}[style=php, caption={\texttt{app/Exceptions/AccountDeactivatedException.php}
  --- Renders as HTTP 403 Forbidden.}]
<?php

namespace App\Exceptions;

use Exception;
use Illuminate\Http\JsonResponse;

class AccountDeactivatedException extends Exception
{
    public function render(): JsonResponse
    {
        return response()->json([
            'success' => false,
            'message' => 'Your account has been deactivated. '
                       . 'Please contact your administrator.',
        ], 403);
    }
}
\end{lstlisting}

%% ─────────────────────────────────────────────────────────────────────────
\section{Security Properties: Summary}
\label{sec:auth-security-summary}
%% ─────────────────────────────────────────────────────────────────────────

The table below provides a consolidated view of the security properties
provided by the authentication and authorisation subsystem.

\begin{table}[h!]
  \centering
  \caption{Security properties of the authentication and authorisation
    subsystem and the mechanism that provides each property.}
  \label{tab:security-properties}
  \renewcommand{\arraystretch}{1.5}
  \begin{tabularx}{\linewidth}{@{}L{4.2cm} X@{}}
    \toprule
    \rowcolor{PrimaryBlue}
    \tblhdr{Security Property} & \tblhdr{Mechanism} \\
    \midrule
    Token integrity
      & RSA-signed Personal Access Tokens issued by Laravel Passport;
        tampering with the token payload invalidates the signature. \\
    \rowalt
    Server-side revocation
      & \texttt{revoked} flag in \texttt{oauth\_access\_tokens};
        logout atomically invalidates the token regardless of its
        remaining cryptographic validity period. \\
    Stateless authentication
      & No server-side session state; the bearer token is the sole proof
        of identity on every request. Enables horizontal scaling. \\
    \rowalt
    Role-based access control
      & \phpclass{RoleMiddleware} enforces non-overlapping, strictly
        bounded permission surfaces at the route-group level. \\
    Clinic data isolation
      & \phpclass{ClinicScopeMiddleware} injects authoritative
        \texttt{clinic\_id}; explicit ownership checks in all
        resource-by-UUID controllers. \\
    \rowalt
    Vertical escalation prevention
      & Role read from authenticated token owner's DB record only;
        self-demotion guard; role middleware applied before controller
        entry. \\
    Horizontal escalation prevention
      & Two-layer clinic scope enforcement (middleware injection +
        controller ownership check); secondary doctor-ownership check
        for prescriptions. \\
    \rowalt
    User enumeration resistance
      & \phpclass{ClinicScopeViolationException} produces an identical
        \stforbid{} response for both ``wrong clinic'' and ``role
        mismatch'' cases; deactivated account message reveals nothing
        about other accounts. \\
    Double-booking prevention
      & Distributed \texttt{Cache::lock()} around the book/reschedule
        write path; unique DB index as final backstop. \\
    \rowalt
    Password confidentiality
      & Passwords stored as Bcrypt hashes via Laravel's built-in
        \texttt{hashed} cast; plaintext is never persisted or logged. \\
    Token confidentiality
      & Raw token returned to client exactly once at issuance; server
        stores only a SHA-256 hash of the secret component in
        \texttt{oauth\_access\_tokens}. \\
    \bottomrule
  \end{tabularx}
\end{table}
